%----------------------------------------------------------------------------------------
% Tailored CV — Doktorand, Bioinformatik / Hepatologie / Onkologie
% Medizinische Klinik 1, Universitätsklinikum Erlangen
% Ref. he_030_PD — PD Dr. Peter Dietrich
%----------------------------------------------------------------------------------------

\documentclass[10pt,a4paper,sans]{moderncv}

\usepackage[utf8]{inputenc}
\usepackage[T1]{fontenc}

\moderncvstyle{classic}
\moderncvcolor{blue}

\usepackage[scale=0.78]{geometry}
\usepackage{ragged2e}

%----------------------------------------------------------------------------------------

\firstname{Kundai Farai}
\familyname{Sachikonye}

\address{Auf der Lichtung 19}{82178 Puchheim, Germany}
\mobile{(+49) 1626386344}
\email{kundai.sachikonye@bitspark.com}
\homepage{github.com/fullscreen-triangle}{github.com/fullscreen-triangle}

%----------------------------------------------------------------------------------------

\begin{document}

\makecvtitle

%----------------------------------------------------------------------------------------
%	EDUCATION
%----------------------------------------------------------------------------------------

\section{Education}

\cventry{2019--2021}{PhD candidate, Computational Lipidomics}{Technische Universität München / Bayerisches Forschungsinstitut}{}{}{
  ML for large-scale MS omics data in clinical cohorts; multi-omics statistical
  modelling; distributed inference across federated clinical sites.
  Departed to pursue independent research.
}
\cventry{2018--2019}{MSc Computational Biology}{Universität Tübingen}{}{}{
  Probabilistic graphical models, statistical bioinformatics, Bayesian methods,
  computational systems biology, ML for biological sequences.
}
\cventry{2014--2017}{MSc Pharmaceutical Biotechnology}{HAW Hamburg}{}{}{
  Molecular biology, cell culture, fermentation, analytical chemistry,
  quality assurance in pharmaceutical production.
}
\cventry{2011--2014}{BSc Biochemistry and Cell Biology}{Jacobs University Bremen}{}{}{
  Biochemistry, cell biology, molecular genetics, protein chemistry.
}

%----------------------------------------------------------------------------------------
%	RELEVANT MANUSCRIPTS
%----------------------------------------------------------------------------------------

\section{Relevant Manuscripts}

\cvitem{2025}{
  \textbf{On the Thermodynamic Consequences of Categorical Completion on Biological Systems.}
  Spectral complementarity framework for immunopeptidology: MHC-I/II binding
  AUC\,$0.81$ vs.\ NetMHCpan\,$0.68$ ($r=0.73$, $p<10^{-12}$, $n=2{,}847$
  IEDB peptides); neoantigen identification from somatic mutations;
  immune escape predictor for tumour immunology.
  \textit{Manuscript. AIMe Registry for Artificial Intelligence.}
}
\cvitem{2025}{
  \textbf{On Disease Epistemology in Blind Receiver.}
  Hierarchical Bayesian disease state $\mathbf{D}\in[0,1]^8$; AUC$=1.00$;
  25/25 validation experiments; sparse $\ell_1$ LP for personalised therapeutic
  target selection from multi-omics data.
  \textit{Manuscript. AIMe Registry / TUM Weihenstephan.}
}

%----------------------------------------------------------------------------------------
%	BIOINFORMATICS AND SEQUENCING TOOLS
%----------------------------------------------------------------------------------------

\section{Bioinformatics and Sequencing Analysis}

\cvitem{gospel}{
  \textbf{RNA-seq and Genomics Analysis Pipeline.}
  Read alignment and quantification; differential expression across conditions
  and cohorts; splicing analysis; somatic and germline variant calling;
  species-aware analysis with ortholog mapping for human/murine comparison;
  cross-species pathway conservation.
  Bayesian network multi-omics integration (genomics $\to$ transcriptomics $\to$
  metagenomics as conditional probability graph).
  $10^6$ variants/second; ClinVar, KEGG, Reactome annotation; gnomAD, 1000~Genomes,
  UK~Biobank validated. Ray distributed, Docker deployable.
  \href{https://github.com/fullscreen-triangle/gospel}{github.com/fullscreen-triangle/gospel}
}

\cvitem{single-cell}{
  \textbf{Single-Cell RNA-seq Statistical Infrastructure.}
  Variational Bayes posterior approximation at $\mathcal{O}(10^6)$ latent variables;
  Gaussian mixture models for rare cell population identification (hepatic stellate
  cells, Kupffer cells, sinusoidal endothelial cells, periportal/pericentral
  hepatocytes); Gaussian process trajectory inference (fibrogenesis progression,
  tumour microenvironment remodelling, immune recruitment).
  Scanpy and Seurat pipelines (UMAP, Leiden, RNA velocity, diffusion pseudotime);
  AnnData / HDF5; probabilistic dimensionality reduction underlying UMAP geometry.
}

\cvitem{syndrome}{
  \textbf{HCC Immunopeptidology Toolkit.}
  Tumour neoantigen identification from somatic mutation data;
  MHC-I/II binding prediction (AUC\,$0.81$); immune escape scoring via
  substitution scan $E(v)=-\max\rho(v,\mathrm{self})+w\cdot\rho(v,\mathrm{target})$;
  disease state $\mathbf{D}\in[0,1]^8$ including $D_\mathrm{Imm}$
  (tumour immune activation) for TME characterisation from RNA-seq data.
  Live: \href{https://syndrome-seven.vercel.app/tools}{syndrome-seven.vercel.app/tools}.
}

\cvitem{lavoisier}{
  \textbf{Mass Spectrometry Omics Pipeline.}
  MS1/MS2 processing + CNN spectral classification; 98.9\,\% NIST accuracy;
  $>100$\,GB mzML at clinical cohort scale.
  Applicable to proteomics and metabolomics integration with RNA-seq data.
  \href{https://github.com/fullscreen-triangle/lavoisier}{github.com/fullscreen-triangle/lavoisier}
}

%----------------------------------------------------------------------------------------
%	EXPERIENCE
%----------------------------------------------------------------------------------------

\section{Experience}

\cventry{2021--present}{Independent Computational Biology Researcher}{Bitspark GmbH}{Munich}{}{
  Computational immunopeptidology, large-scale genomics, Bayesian multi-omics
  integration, and clinical data infrastructure for oncology and pharmacology.
}

\cventry{2019--2021}{Computational Lipidomics}{TUM / Bayerisches Forschungsinstitut}{Munich}{}{
  ML for clinical MS omics data; distributed pipeline maintenance; multi-omics
  integration across federated clinical sites. Close collaboration with wet lab
  teams: understanding sample preparation constraints and protocol artefacts
  in downstream data. Python, Java.
}

\cventry{2017--2018}{Glycomics and Proteomics}{Universitätsklinikum Hamburg-Eppendorf}{}{}{
  Wet lab: automated LC-MS/MS and MALDI-TOF pipeline operation;
  clinical sample processing (tissue, serum, CSF); data annotation and
  quality control for multi-centre clinical cohort studies.
}

%----------------------------------------------------------------------------------------
%	TECHNICAL SKILLS
%----------------------------------------------------------------------------------------

\section{Technical Skills}

\cvitem{Sequencing Analysis}{
  RNA-seq (alignment, quantification, DE, splicing); scRNA-seq (Scanpy, Seurat,
  UMAP, Leiden, RNA velocity, diffusion pseudotime); human/murine comparative
  transcriptomics; FASTQ, BAM, VCF; samtools, Nextflow; AnnData / HDF5
}

\cvitem{Statistical / ML}{
  R (Bioconductor, DESeq2, edgeR, limma, ggplot2); Python (PyTorch, scikit-learn,
  scanpy, pandas); variational Bayes; probabilistic graphical models; Gaussian processes;
  MCMC (NUTS/HMC); sparse regularisation ($\ell_1$); GMMs
}

\cvitem{Oncology / Immunology}{
  Neoantigen identification; MHC binding prediction; tumour immune escape scoring;
  tumour microenvironment characterisation from scRNA-seq; disease state
  quantification ($\mathbf{D}\in[0,1]^8$); therapeutic target prioritisation
}

\cvitem{Wet Lab}{
  LC-MS/MS, MALDI-TOF; cell culture; clinical sample processing (tissue, serum);
  PCR; protein biochemistry; pharmaceutical production (fermentation, downstream)
}

%----------------------------------------------------------------------------------------
%	LANGUAGES
%----------------------------------------------------------------------------------------

\section{Languages}

\cvitemwithcomment{English}{Native / Fluent}{}
\cvitemwithcomment{German}{Conversational}{Living in Germany continuously since 2011}

%----------------------------------------------------------------------------------------

\renewcommand{\listitemsymbol}{-~}

\end{document}
